\section{Union Closure for Normal PEPS}\label{sec:peps_normal_union}

A normal PEPS is injective only after blocking certain regions, so the injective
theorem applies once those regions are available. The closure property that
makes them available is that, for positive bond dimensions, the union of two
injective regions is again injective,
$\inj(A)\wedge\inj(B)\Longrightarrow\inj(A\cup B)$.
Writing $P_0=A\setminus B$, $P_1=A\cap B$, $P_2=B\setminus A$, and
$P_3=(A\cup B)^c$, a boundary operator $X$ attached to $P_3$ that is annihilated
by the contraction over $A\cup B$ descends through the contraction chain
\begin{align}
    \mathcal C_{\{0,1,2\}}(X)=0
    \Longrightarrow \mathcal C_{\{2\}}(X)=0
    \Longrightarrow \mathcal C_{\{1,2\}}(X)=0
    \Longrightarrow X=0,
    \notag
\end{align}
so the union contraction has trivial kernel. Iterating gives injectivity of
every finite union of injective regions.

\begin{definition}[Four-region decomposition for a union of regions]%
    \label{def:peps_regionUnionDecomposition}
    \lean{TNLean.PEPS.regionOnlyLeft, TNLean.PEPS.regionOverlap}
    \lean{TNLean.PEPS.regionOnlyRight, TNLean.PEPS.regionOutsideUnion}
    \lean{TNLean.PEPS.regionUnionPart}
    \leanok
    Let $V$ be a finite vertex set. For two regions $A,B\subseteq V$,
    define the four regions $A\setminus B$, $A\cap B$,
    $B\setminus A$, and $V\setminus(A\cup B)$. These are the four
    regions used in the proof of Lemma~\ref{lem:peps_injectiveRegion_union}.
    They also form an indexed family over $\{0,1,2,3\}$ in the order
    $A\setminus B$, $A\cap B$, $B\setminus A$, $V\setminus(A\cup B)$.
\end{definition}

\begin{lemma}[Identities for the four-region decomposition]%
    \label{lem:peps_regionUnionDecomposition_identities}
    \lean{TNLean.PEPS.regionOnlyLeft_union_overlap}
    \lean{TNLean.PEPS.regionOverlap_union_onlyRight}
    \lean{TNLean.PEPS.regionOutsideUnion_eq_regionComplement_union}
    \lean{TNLean.PEPS.regionThreePart_union, TNLean.PEPS.regionFourPart_union}
    \lean{TNLean.PEPS.regionFourPart_cases}
    \lean{TNLean.PEPS.regionUnionPart_zero, TNLean.PEPS.regionUnionPart_one}
    \lean{TNLean.PEPS.regionUnionPart_two, TNLean.PEPS.regionUnionPart_three}
    \lean{TNLean.PEPS.regionUnionPart_three_eq_regionComplement_union}
    \lean{TNLean.PEPS.regionUnionPart_zero_union_one}
    \lean{TNLean.PEPS.regionUnionPart_biUnion_zero_one}
    \lean{TNLean.PEPS.regionUnionPart_one_union_two}
    \lean{TNLean.PEPS.regionUnionPart_biUnion_one_two}
    \lean{TNLean.PEPS.regionUnionPart_zero_union_one_union_two}
    \lean{TNLean.PEPS.regionUnionPart_biUnion_zero_one_two}
    \lean{TNLean.PEPS.regionUnionPart_biUnion_three}
    \lean{TNLean.PEPS.regionUnionPart_biUnion_three_eq_regionComplement_zero_one_two}
    \lean{TNLean.PEPS.regionUnionPart_three_eq_regionComplement_left_inter_right}
    \lean{TNLean.PEPS.regionUnionPart_two_union_three_eq_regionComplement_left}
    \lean{TNLean.PEPS.regionUnionPart_biUnion_two_three_eq_regionComplement_left}
    \lean{TNLean.PEPS.regionUnionPart_biUnion_two_three_eq_regionComplement_zero_one}
    \lean{TNLean.PEPS.regionUnionPart_zero_union_three_eq_regionComplement_right}
    \lean{TNLean.PEPS.regionUnionPart_biUnion_zero_three_eq_regionComplement_right}
    \lean{TNLean.PEPS.regionUnionPart_biUnion_zero_three_eq_regionComplement_one_two}
    \lean{TNLean.PEPS.regionUnionPart_two_union_three_eq_regionComplement_zero_union_one}
    \lean{TNLean.PEPS.regionUnionPart_zero_union_three_eq_regionComplement_one_union_two}
    \lean{TNLean.PEPS.regionUnionPart_zero_union_one_eq_regionComplement_two_union_three}
    \lean{TNLean.PEPS.regionUnionPart_one_union_two_eq_regionComplement_zero_union_three}
    \lean{TNLean.PEPS.regionUnionPart_biUnion}
    \lean{TNLean.PEPS.regionUnionPart_exists_unique}
    \leanok
    \uses{def:peps_regionUnionDecomposition}
    Let $V$ be a finite vertex set and let $A,B\subseteq V$. The left-only
    and overlap regions reconstruct $A$, the overlap and right-only regions
    reconstruct $B$, the three inside regions reconstruct $A\cup B$, and
    all four regions reconstruct $V$. The same identities also hold in the
    indexed four-region family: pieces $0$ and $1$ reconstruct $A$, pieces
    $1$ and $2$ reconstruct $B$, and pieces $0,1,2$ reconstruct $A\cup B$.
    These same reconstructions are available as finite unions over the
    corresponding subfamilies of the indexed four-region family.
    Piece $3$ is the complement of $A\cup B$, equivalently the complement of
    the three inside indexed pieces, and also $A^c\cap B^c$. The singleton
    subfamily $\{3\}$ gives the outside block. Pieces $2,3$
    reconstruct $A^c$, equivalently the complement of pieces $0,1$, and
    pieces $0,3$ reconstruct $B^c$, equivalently the complement of pieces
    $1,2$, again both as explicit unions and as finite unions over the
    indexed subfamilies. Conversely, pieces $0,1$ are the
    complement of pieces $2,3$,
    pieces $1,2$ are the complement of pieces $0,3$, and adding all four
    pieces reconstructs $V$.
    Equivalently, each vertex lies in one of the four regions, and the indexed
    four-region family has union $V$. In fact, every vertex belongs to exactly
    one indexed piece.
\end{lemma}
\begin{proof}\leanok
    Check membership of an arbitrary vertex in $A$ and in $B$.
\end{proof}

\begin{lemma}[Disjointness in the four-region decomposition]%
    \label{lem:peps_regionUnionDecomposition_disjoint}
    \lean{TNLean.PEPS.regionOnlyLeft_disjoint_overlap}
    \lean{TNLean.PEPS.regionOnlyRight_disjoint_overlap}
    \lean{TNLean.PEPS.regionOnlyLeft_disjoint_onlyRight}
    \lean{TNLean.PEPS.regionOnlyLeft_disjoint_outside}
    \lean{TNLean.PEPS.regionOverlap_disjoint_outside}
    \lean{TNLean.PEPS.regionOnlyRight_disjoint_outside}
    \lean{TNLean.PEPS.regionInside_disjoint_outside}
    \lean{TNLean.PEPS.regionUnionPart_zero_one_disjoint_two_three}
    \lean{TNLean.PEPS.regionUnionPart_biUnion_zero_one_disjoint_two_three}
    \lean{TNLean.PEPS.regionUnionPart_one_two_disjoint_zero_three}
    \lean{TNLean.PEPS.regionUnionPart_biUnion_one_two_disjoint_zero_three}
    \lean{TNLean.PEPS.regionUnionPart_inside_disjoint_three}
    \lean{TNLean.PEPS.regionUnionPart_pairwise_disjoint}
    \leanok
    \uses{def:peps_regionUnionDecomposition,
        lem:peps_regionUnionDecomposition_identities}
    In the four-region decomposition, the left-only and right-only parts are
    disjoint from the overlap and from each other. Each inside part, and hence
    the union of the three inside regions, is disjoint from the outside region.
    In indexed form, the pieces forming $A$ are disjoint from the pieces
    forming $A^c$, and the pieces forming $B$ are disjoint from the pieces
    forming $B^c$, both as explicit binary unions and as finite unions over
    the corresponding selected indexed subfamilies. The indexed four-region
    family is pairwise disjoint.
\end{lemma}
\begin{proof}\leanok
    Use the standard disjointness identities for set difference and
    intersection, together with the reconstruction of $A\cup B$.
\end{proof}

\begin{definition}[Contraction maps for the union injectivity proof]%
    \label{def:peps_injectiveUnionContractionMaps}
    \lean{TNLean.PEPS.InjectiveUnionContractionMaps}
    \leanok
    Let $K$, $E_{012}$, $E_2$, and $E_{12}$ be complex vector spaces. A
    contraction-map family for the union-injectivity lemma in
    \cite[Section~4]{Molnar2018NormalPEPS} consists of three linear maps
    $\mathcal C_{\{0,1,2\}}:K\to E_{012}$,
    $\mathcal C_{\{2\}}:K\to E_2$, and
    $\mathcal C_{\{1,2\}}:K\to E_{12}$.
\end{definition}

\begin{definition}[Contraction chain for the union injectivity proof]%
    \label{def:peps_injectiveUnionContractionChain}
    \lean{TNLean.PEPS.InjectiveUnionContractionChain}
    \leanok
    \uses{def:peps_injectiveUnionContractionMaps}
    A contraction chain consists of the three contraction maps of
    Definition~\ref{def:peps_injectiveUnionContractionMaps}, together with the
    three implications used in
    Lemma~\ref{lem:peps_injectiveRegion_union_contractionChain}.
\end{definition}

\begin{lemma}[Kernel criterion for the union contraction chain]%
    \label{lem:peps_injectiveRegion_union_contractionChain}
    \lean{TNLean.PEPS.InjectiveUnionContractionChain.union_contraction_ker_eq_bot}
    \lean{TNLean.PEPS.InjectiveUnionContractionChain.union_contraction_injective}
    \leanok
    \uses{def:peps_injectiveUnionContractionChain}
    Let $P_0=A\setminus B$, $P_1=A\cap B$, $P_2=B\setminus A$, and
    $P_3=(A\cup B)^c$. Let $K$ be the boundary vector space attached to
    $P_3$, and let
    $\mathcal C_{\{0,1,2\}}:K\to E_{012}$,
    $\mathcal C_{\{2\}}:K\to E_2$, and
    $\mathcal C_{\{1,2\}}:K\to E_{12}$ be the three contraction maps in the
    proof of the union-injectivity lemma in
    \cite[Section~4]{Molnar2018NormalPEPS}. If
    \begin{align}
        \mathcal C_{\{0,1,2\}}(X)=0
        &\Longrightarrow \mathcal C_{\{2\}}(X)=0,
        \notag\\
        \mathcal C_{\{2\}}(X)=0
        &\Longrightarrow \mathcal C_{\{1,2\}}(X)=0,
        \notag
    \end{align}
    and $\mathcal C_{\{1,2\}}(X)=0\Longrightarrow X=0$, then
    $\ker \mathcal C_{\{0,1,2\}}=\{0\}$.
\end{lemma}
\begin{proof}\leanok
    For $X\in K$ with $\mathcal C_{\{0,1,2\}}(X)=0$, apply the three
    implications in sequence:
    \begin{align}
        \mathcal C_{\{0,1,2\}}(X)=0
        \Longrightarrow \mathcal C_{\{2\}}(X)=0
        \Longrightarrow \mathcal C_{\{1,2\}}(X)=0
        \Longrightarrow X=0.
        \notag
    \end{align}
\end{proof}

\begin{lemma}[Union of injective regions]%
    \label{lem:peps_injectiveRegion_union}
    \notready
    \uses{def:IsVertexInjective, lem:peps_regionUnionDecomposition_identities,
        lem:peps_regionUnionDecomposition_disjoint,
        lem:peps_injectiveRegion_union_contractionChain}
    Let $V$ be a finite vertex set. Write $\inj(R)$ for the
    assertion that the PEPS tensor blocked to the region $R\subseteq V$ is
    injective. For two possibly overlapping regions $A,B\subseteq V$,
    $\inj(A)\wedge\inj(B)\Longrightarrow\inj(A\cup B)$.
    The union-injectivity proof in \cite[Section~4]{Molnar2018NormalPEPS}
    blocks the network into four parts $A\setminus B$, $A\cap B$, $B\setminus A$, and
    $(A\cup B)^c$, and applies the two available left inverses in sequence:
    % Formula: the four blocked tensors are P_0=A\setminus B, P_1=A\cap B,
    %   P_2=B\setminus A, and P_3=(A\cup B)^c.
    % Source verdict: Molnar et al., arXiv:1804.04964, page 14, lines
    %   1329--1345 of paper_normal.tex.  The source draws the complete graph
    %   on the four blocked tensors and one upper physical leg at each tensor.
    % Ink-to-index map: every one of the six joins between distinct tensor
    %   blocks contracts one virtual index; the four upper joins end at
    %   inkless boundary atoms and carry the four free physical indices.
    % Type and contraction verdict: all six internal joins are
    %   virtual--virtual, all four open-leg joins are physical--physical, and
    %   the internal virtual graph is exactly K_4.
    % Structural boundary verdict: the typed boundary atoms and their joins
    %   derive (V,P)=(0,4).  The free tier does not emit a boundary event.
    \begin{center}
        \begin{tenkzfree}[compact]
            \tnput[dot, label pos=west,
                ports={45:virtual,east:virtual,-45:virtual,north:physical}]
                {unionLeft}{(-14mm,0)}{A\setminus B}
            \tnput[dot, label pos=east,
                ports={-135:virtual,-35:virtual,-70:virtual,north:physical}]
                {unionOverlap}{(-4mm,11mm)}{A\cap B}
            \tnput[dot, label pos=east,
                ports={west:virtual,135:virtual,-135:virtual,north:physical}]
                {unionRight}{(14mm,0)}{B\setminus A}
            \tnput[dot, label pos=south,
                ports={150:virtual,110:virtual,45:virtual,north:physical}]
                {unionOutside}{(4mm,-11mm)}{(A\cup B)^c}

            \tnput[boundary, ports={south:physical}]
                {unionLeftPhysical}{(-14mm,7mm)}{}
            \tnput[boundary, ports={south:physical}]
                {unionOverlapPhysical}{(-4mm,18mm)}{}
            \tnput[boundary, ports={south:physical}]
                {unionRightPhysical}{(14mm,7mm)}{}
            \tnput[boundary, ports={south:physical}]
                {unionOutsidePhysical}{(4mm,-4mm)}{}

            \tnjoin{unionLeft.45}{unionOverlap.-135}
            \tnjoin{unionLeft.east}{unionRight.west}
            \tnjoin{unionLeft.-45}{unionOutside.150}
            \tnjoin{unionOverlap.-35}{unionRight.135}
            \tnjoin{unionOverlap.-70}{unionOutside.110}
            \tnjoin{unionRight.-135}{unionOutside.45}
            \tnjoin{unionLeft.north}{unionLeftPhysical.south}
            \tnjoin{unionOverlap.north}{unionOverlapPhysical.south}
            \tnjoin{unionRight.north}{unionRightPhysical.south}
            \tnjoin{unionOutside.north}{unionOutsidePhysical.south}
        \end{tenkzfree}
    \end{center}
    Taking a boundary tensor $X$ on $(A\cup B)^c$ with
    $\mathcal C_{\{0,1,2\}}(X)=0$, injectivity of $A$ gives
    $\mathcal C_{\{2\}}(X)=0$; reinserting $A\cap B$ and applying
    injectivity of $B$ then gives $X=0$:
    % Formula: C_{\{0,1,2\}}(X)=0 -> C_{\{2\}}(X)=0
    %   -> C_{\{1,2\}}(X)=0 -> X=0.
    % Source verdict: Molnar et al., arXiv:1804.04964, page 14, lines
    %   1346--1401 of paper_normal.tex.  The arrows are, in order, the left
    %   inverse on A, reinsertion of A\cap B, and the left inverse on B.
    % Ink-to-index map: panel 1 is K_4 on P_0,P_1,P_2,X.  Panel 2 retains the
    %   P_2--X contraction and exposes their four bonds to the removed
    %   P_0,P_1 block.  Panel 3 restores P_1 and its triangle with P_2,X,
    %   leaving the three bonds to P_0 open.  Panel 4 leaves the three virtual
    %   indices incident on X open.  Upper legs are free physical indices.
    % Type and contraction verdict: every tensor--tensor join is
    %   virtual--virtual; every open stub terminates at a typed inkless
    %   boundary atom.  The panels have respectively 6, 1, 3, and 0 internal
    %   virtual contractions.
    % Structural boundary verdict: typed boundary atoms and their joins derive,
    %   from left to right, (V,P)=(0,3), (4,1), (3,2), and (3,0).  The free
    %   tier does not emit boundary events.
    \begin{center}
        \begin{tenkzfree}[compact]
            \tnput[dot, label pos=west,
                ports={45:virtual,east:virtual,-45:virtual,north:physical}]
                {proofZeroLeft}{(-7mm,0)}{P_0}
            \tnput[dot, label pos=east,
                ports={-135:virtual,-45:virtual,-90:virtual,north:physical}]
                {proofZeroOverlap}{(0,6mm)}{P_1}
            \tnput[dot, label pos=east,
                ports={west:virtual,135:virtual,-135:virtual,north:physical}]
                {proofZeroRight}{(7mm,0)}{P_2}
            \tnput[dot, label pos=south,
                ports={135:virtual,90:virtual,45:virtual}]
                {proofZeroX}{(0,-6mm)}{X}
            \tnput[boundary, ports={south:physical}]
                {proofZeroLeftPhysical}{(-7mm,4mm)}{}
            \tnput[boundary, ports={south:physical}]
                {proofZeroOverlapPhysical}{(0,10mm)}{}
            \tnput[boundary, ports={south:physical}]
                {proofZeroRightPhysical}{(7mm,4mm)}{}

            \tnjoin{proofZeroLeft.45}{proofZeroOverlap.-135}
            \tnjoin{proofZeroLeft.east}{proofZeroRight.west}
            \tnjoin{proofZeroLeft.-45}{proofZeroX.135}
            \tnjoin{proofZeroOverlap.-45}{proofZeroRight.135}
            \tnjoin{proofZeroOverlap.-90}{proofZeroX.90}
            \tnjoin{proofZeroRight.-135}{proofZeroX.45}
            \tnjoin{proofZeroLeft.north}{proofZeroLeftPhysical.south}
            \tnjoin{proofZeroOverlap.north}{proofZeroOverlapPhysical.south}
            \tnjoin{proofZeroRight.north}{proofZeroRightPhysical.south}
        \end{tenkzfree}
        \,$=0$\quad$\xrightarrow{A^{-1}}$\quad
        \begin{tenkzfree}[compact]
            \tnput[dot, label pos=east,
                ports={-135:virtual,135:virtual,-90:virtual,north:physical}]
                {proofOneRight}{(4mm,1mm)}{P_2}
            \tnput[dot, label pos=south,
                ports={45:virtual,90:virtual,150:virtual}]
                {proofOneX}{(-2mm,-6mm)}{X}
            \tnput[boundary, ports={east:virtual}]
                {proofOneRightLower}{(-3mm,-1mm)}{}
            \tnput[boundary, ports={-45:virtual}]
                {proofOneRightUpper}{(-2mm,7mm)}{}
            \tnput[boundary, ports={east:virtual}]
                {proofOneXLower}{(-8mm,-3mm)}{}
            \tnput[boundary, ports={south:virtual}]
                {proofOneXUpper}{(-2mm,1mm)}{}
            \tnput[boundary, ports={south:physical}]
                {proofOneRightPhysical}{(4mm,7mm)}{}

            \tnjoin{proofOneRight.-90}{proofOneX.45}
            \tnjoin{proofOneRight.-135}{proofOneRightLower.east}
            \tnjoin{proofOneRight.135}{proofOneRightUpper.-45}
            \tnjoin{proofOneX.150}{proofOneXLower.east}
            \tnjoin{proofOneX.90}{proofOneXUpper.south}
            \tnjoin{proofOneRight.north}{proofOneRightPhysical.south}
        \end{tenkzfree}
        \,$=0$\quad$\xrightarrow{A\cap B}$\quad
        \begin{tenkzfree}[compact]
            \tnput[dot, label pos=east,
                ports={-135:virtual,-90:virtual,-45:virtual,north:physical}]
                {proofTwoOverlap}{(0,6mm)}{P_1}
            \tnput[dot, label pos=east,
                ports={135:virtual,-135:virtual,west:virtual,north:physical}]
                {proofTwoRight}{(7mm,0)}{P_2}
            \tnput[dot, label pos=south,
                ports={90:virtual,45:virtual,135:virtual}]
                {proofTwoX}{(0,-6mm)}{X}
            \tnput[boundary, ports={45:virtual}]
                {proofTwoOverlapFree}{(-5mm,2mm)}{}
            \tnput[boundary, ports={east:virtual}]
                {proofTwoRightFree}{(-5mm,0)}{}
            \tnput[boundary, ports={-45:virtual}]
                {proofTwoXFree}{(-5mm,-2mm)}{}
            \tnput[boundary, ports={south:physical}]
                {proofTwoOverlapPhysical}{(0,10mm)}{}
            \tnput[boundary, ports={south:physical}]
                {proofTwoRightPhysical}{(7mm,4mm)}{}

            \tnjoin{proofTwoOverlap.-135}{proofTwoOverlapFree.45}
            \tnjoin{proofTwoRight.west}{proofTwoRightFree.east}
            \tnjoin{proofTwoX.135}{proofTwoXFree.-45}
            \tnjoin{proofTwoOverlap.-90}{proofTwoX.90}
            \tnjoin{proofTwoOverlap.-45}{proofTwoRight.135}
            \tnjoin{proofTwoRight.-135}{proofTwoX.45}
            \tnjoin{proofTwoOverlap.north}{proofTwoOverlapPhysical.south}
            \tnjoin{proofTwoRight.north}{proofTwoRightPhysical.south}
        \end{tenkzfree}
        \,$=0$\quad$\xrightarrow{B^{-1}}$\quad
        \begin{tenkzfree}[compact]
            \tnput[dot, label pos=south, ports={135:virtual,90:virtual,45:virtual}]
                {proofThreeX}{(0,0)}{X}
            \tnput[boundary, ports={-45:virtual}]
                {proofThreeLeft}{(-4mm,5mm)}{}
            \tnput[boundary, ports={south:virtual}]
                {proofThreeMiddle}{(0,6mm)}{}
            \tnput[boundary, ports={-135:virtual}]
                {proofThreeRight}{(4mm,5mm)}{}

            \tnjoin{proofThreeX.135}{proofThreeLeft.-45}
            \tnjoin{proofThreeX.90}{proofThreeMiddle.south}
            \tnjoin{proofThreeX.45}{proofThreeRight.-135}
        \end{tenkzfree}
        \,$=0$.
    \end{center}
\end{lemma}
\begin{proof}
    Put $P_0=A\setminus B$, $P_1=A\cap B$, $P_2=B\setminus A$, and
    $P_3=(A\cup B)^c$. Then $P_0\cup P_1=A$, $P_1\cup P_2=B$, and
    $P_0\cup P_1\cup P_2=A\cup B$. Let $\mathcal C_I$ denote contraction of
    the blocked tensor over the region $\bigcup_{i\in I}P_i$. To prove
    injectivity of $A\cup B$, take a boundary tensor $X$ on $P_3$ such that
    $\mathcal C_{\{0,1,2\}}(X)=0$. Since $P_0\cup P_1=A$ is injective, its
    left inverse gives
    $\mathcal C_{\{0,1,2\}}(X)=0\Longrightarrow
    \mathcal C_{\{2\}}(X)=0$.
    The remaining two steps are
    $\mathcal C_{\{2\}}(X)=0\Longrightarrow
    \mathcal C_{\{1,2\}}(X)=0$ and
    $\mathcal C_{\{1,2\}}(X)=0\Longrightarrow X=0$.
    The first implication is obtained by contracting with the tensor on
    $P_1$, and the second is the left inverse for the injective region
    $P_1\cup P_2=B$.
    Thus $\ker \mathcal C_{\{0,1,2\}}=\{0\}$, and therefore
    $P_0\cup P_1\cup P_2=A\cup B$ is injective.
\end{proof}

\begin{lemma}[Union of the edge-centred blocks]%
    \label{lem:peps_injectiveRegion_union_edgeBlocked}
    \lean{TNLean.PEPS.regionBlockedTensorInjective_union}
    \leanok
    \uses{def:peps_regionInjectivityDataOf}
    Fix an edge $e$ and a partition of the vertex set into the three
    edge-centred blocks $R$, $B$, $C$ (red, blue, complementary) of the normal
    PEPS proof, with the red block carrying one endpoint of $e$ and the blue
    block the other. The injectivity of the blue and complementary blocks is
    part of the edge-centred blocking data; if every virtual bond dimension
    is positive, then their union $B\cup C=V\setminus R$ is injective.
\end{lemma}
\begin{proof}\leanok
    This is Lemma~\ref{lem:peps_injectiveRegion_union} for the
    edge-centred triple, where the two regions are the blue and complementary
    blocks and their union is the set complement of the red block. A coefficient
    family annihilating the blocked weights of $V\setminus R$ is stripped of the
    blue block by its left inverse, leaving a complement-coupling combination
    that vanishes for every blue boundary configuration. The complementary block
    then carries each coupling coefficient, read as a function of the
    complementary physical leg, into a combination of its blocked weights whose
    coefficient at a boundary configuration is the host residual coefficient
    reconstructed from the blue and complementary boundary data; injectivity of
    the complementary block forces that coefficient to vanish wherever a global
    configuration realizes the two boundary labels. Since every host residual is
    realized once the bond dimensions are positive, the coefficient family is
    zero.
\end{proof}

\begin{lemma}[Union of two disjoint injective regions]%
    \label{lem:peps_injectiveRegion_union_disjoint}
    \lean{TNLean.PEPS.regionBlockedTensorInjective_union_disjoint}
    \leanok
    \uses{lem:peps_injectiveRegion_union_edgeBlocked}
    Let $S,T\subseteq V$ be two disjoint regions. If the blocked tensors of $S$
    and of $T$ are each injective and every virtual bond dimension is positive,
    then the blocked tensor of their union $S\cup T$ is injective.
\end{lemma}
\begin{proof}\leanok
    This is Lemma~\ref{lem:peps_injectiveRegion_union} in the case
    where the two regions are disjoint, so that the overlap $S\cap T$ is empty
    and the four-region picture collapses to the three blocks $S$, $T$, and the
    set complement $(S\cup T)^c$. The argument places $S$ and $T$ as the blue and
    complementary blocks of a three-block decomposition whose third block is the
    set complement of their union, and applies the two-step inverse argument of
    the edge-centred Lemma~\ref{lem:peps_injectiveRegion_union_edgeBlocked}
    to that decomposition. No injectivity hypothesis on the set complement and no
    distinguished edge are needed.
\end{proof}

\begin{lemma}[Union of two injective regions with positive bonds]%
    \label{lem:peps_injectiveRegion_union_pos}
    \lean{TNLean.PEPS.regionBlockedTensorInjective_union_overlap}
    \leanok
    \uses{lem:peps_injectiveRegion_union_edgeBlocked}
    Let $A,B\subseteq V$ be two possibly overlapping regions. If the blocked
    tensors of $A$ and of $B$ are each injective and every virtual bond dimension
    is positive, then the blocked tensor of their union $A\cup B$ is injective.
    This is Lemma~\ref{lem:peps_injectiveRegion_union} with the explicit
    positive-bond hypothesis carried by the formalization; the overlap is no longer
    required to be empty.
\end{lemma}
\begin{proof}\leanok
    Write the four parts $P_0=A\setminus B$, $P_1=A\cap B$, $P_2=B\setminus A$.
    A coefficient family annihilating the blocked weights of $A\cup B$ is stripped
    of $A$ by its left inverse, leaving the coupling through $P_2$ vanishing for
    every $A$-boundary configuration. The source then reinserts the overlap $P_1$
    and inverts $B$. The legs running from $P_0$ to the outside are the open legs
    of the final tensor; they stay uncontracted through the reinsertion. Carrying
    those legs as a separate label, fix it to a reference value and restrict the
    coefficient family to the matching fiber. The fiber-restricted first strip
    still vanishes, so reinserting $P_1$ and inverting $B$ forces the fiber row to
    vanish. The union boundary edges partition into the $B$-boundary edges and the
    $P_0$-outer edges, so fixing both the reference label and the $B$-boundary
    label of a configuration realizing a host label selects that single host term,
    forcing the coefficient to vanish at every realizable host label. Surjectivity
    of the union host boundary, available once the bond dimensions are positive,
    realizes every label, so the coefficient family is zero.
\end{proof}

\begin{lemma}[Finite disjoint union of injective regions]%
    \label{lem:peps_injectiveRegion_biUnion_disjoint}
    \lean{TNLean.PEPS.regionBlockedTensorInjective_biUnion_disjoint}
    \leanok
    \uses{lem:peps_injectiveRegion_union_disjoint}
    Let $(R_i)_{i\in s}$ be a nonempty finite family of pairwise-disjoint regions.
    If the blocked tensor of each $R_i$ is injective and every virtual bond
    dimension is positive, then the blocked tensor of their union
    $\bigcup_{i\in s} R_i$ is injective.
\end{lemma}
\begin{proof}\leanok
    Induct on the nonempty family. The singleton case is the injectivity of the
    one block. For the inductive step, the next block is disjoint from the union
    of the remaining blocks, so
    Lemma~\ref{lem:peps_injectiveRegion_union_disjoint} combines the two
    injective pieces.
\end{proof}

\begin{lemma}[Finite union of injective regions]%
    \label{lem:peps_injectiveRegion_biUnion}
    \lean{TNLean.PEPS.regionBlockedTensorInjective_biUnion_overlap}
    \leanok
    \uses{lem:peps_injectiveRegion_union_pos}
    Let $(R_i)_{i\in s}$ be a nonempty finite family of possibly overlapping
    regions. If the blocked tensor of each $R_i$ is injective and every virtual
    bond dimension is positive, then the blocked tensor of their union
    $\bigcup_{i\in s} R_i$ is injective. This is the finite iteration of the
    overlapping union lemma, with no disjointness assumption.
\end{lemma}
\begin{proof}\leanok
    Induct on the nonempty family. The singleton case is the injectivity of the
    one block. For the inductive step, combine the next block with the union of
    the remaining blocks by Lemma~\ref{lem:peps_injectiveRegion_union_pos}; no
    disjointness is needed.
\end{proof}

\begin{definition}[Injectivity predicate for blocked regions]%
    \label{def:peps_regionInjectivityData}
    \lean{TNLean.PEPS.RegionInjectivityData}
    \lean{TNLean.PEPS.RegionInjectivityUnionClosure}
    \leanok
    A region-injectivity hypothesis assigns to each finite vertex region the
    assertion that the tensor obtained by blocking that region is injective.
    This predicate is used for regions such as $R$, $S$, $T$, and their
    complements in the normal PEPS proof. The union-closure hypothesis records
    the conclusion supplied by the source union-of-injective-regions lemma:
    the union of two injective regions is injective.
\end{definition}

\begin{definition}[Concrete region-injectivity predicate of a tensor]%
    \label{def:peps_regionInjectivityDataOf}
    \lean{TNLean.PEPS.regionInjectivityDataOf}
    \leanok
    \uses{def:peps_regionInjectivityData}
    The region-injectivity predicate of a fixed PEPS tensor $A$ sends a finite
    region $R$ to the assertion that the blocked-region tensor of $R$ is
    injective.
\end{definition}

\begin{theorem}[Gauged blocked-region weight as a global sum]
    \label{thm:peps_regionBlockedWeight_applyGauge_eq_globalSum}
    \lean{TNLean.PEPS.regionBlockedWeight_applyGauge_eq_globalSum}
    \leanok
    \uses{def:applyGauge}
    Let $\widetilde B=\operatorname{Gauge}_X(B)$ be obtained from $B$ by a gauge
    family $X$. For a region $R$, a boundary configuration $\mu$, and a
    physical configuration $\tau$ on $R$, the blocked weight of $\widetilde B$
    is
    \begin{align}
        T_{\widetilde B,R}^{\mu}(\tau)
        &=
        \sum_{\zeta}
        \left(\prod_{f\in\partial R}G_f^X(\mu_f,\zeta_f)\right)
        \prod_{v\in R}B_v(\zeta|_{\inc(v)},\tau_v),
        \label{eq:peps_gauged_global_sum}
    \end{align}
    where the sum ranges over global virtual configurations $\zeta$, and
    $G_f^X$ is the surviving gauge on the boundary edge $f$.
\end{theorem}
\begin{proof}\leanok
    In the contraction defining $T_{\widetilde B,R}^{\mu}(\tau)$, the two
    half-edge gauge factors on every edge internal to $R$ combine as
    \begin{align}
        \sum_a (X_e)_{\alpha,a}(X_e^{-1})_{a,\beta}
        &=
        (X_eX_e^{-1})_{\alpha,\beta}
        =
        \delta_{\alpha,\beta}.
        \notag
    \end{align}
    Thus every interior edge contributes the identity and its gauge cancels.
    The only remaining factors are the boundary gauges
    $G_f^X(\mu_f,\zeta_f)$, giving (\ref{eq:peps_gauged_global_sum}).
\end{proof}

\begin{theorem}[Boundary-coupling form of the gauged blocked weight]
    \label{thm:peps_regionBlockedWeight_applyGauge}
    \lean{TNLean.PEPS.regionBlockedWeight_applyGauge}
    \leanok
    \uses{thm:peps_regionBlockedWeight_applyGauge_eq_globalSum}
    With the same notation, set
    $K_R^X(\mu,\nu):=\prod_{f\in\partial R}G_f^X(\mu_f,\nu_f)$.
    Then the gauged blocked weight is the original blocked weight multiplied by
    the boundary-coupling matrix:
    \begin{align}
        T_{\widetilde B,R}^{\mu}(\tau)
        &=
        \sum_{\nu}K_R^X(\mu,\nu)\,T_{B,R}^{\nu}(\tau).
        \label{eq:peps_boundary_coupling}
    \end{align}
    This is the gauge absorption of the blocked tensor in
    \cite[Section~4, proof of Theorem~3]{Molnar2018NormalPEPS}.
\end{theorem}
\begin{proof}\leanok
    Start from Theorem~\ref{thm:peps_regionBlockedWeight_applyGauge_eq_globalSum}
    and group the global virtual configurations by their boundary label $\nu$:
    \begin{align}
        \sum_{\zeta}
        K_R^X(\mu,\zeta|_{\partial R})
        \prod_{v\in R}B_v(\zeta|_{\inc(v)},\tau_v)
        &=
        \sum_{\nu}K_R^X(\mu,\nu)\,T_{B,R}^{\nu}(\tau).
        \notag
    \end{align}
\end{proof}

\begin{theorem}[The boundary coupling has a right inverse]
    \label{thm:peps_sum_regionBoundaryCoupling_mul_inv}
    \lean{TNLean.PEPS.sum_regionBoundaryCoupling_mul_inv}
    \leanok
    Let
    $K_R^X(\mu,\nu)=\prod_{f\in\partial R}G_f^X(\mu_f,\nu_f)$ and
    $(K_R^X)^{-1}(\nu,\rho)
    =\prod_{f\in\partial R}(G_f^X)^{-1}(\nu_f,\rho_f)$.
    Then, for boundary configurations $\mu$ and $\rho$,
    \begin{align}
        \sum_{\nu}K_R^X(\mu,\nu)\,(K_R^X)^{-1}(\nu,\rho)
        &=
        \begin{cases}
            1, & \mu=\rho,\\
            0, & \mu\neq\rho.
        \end{cases}
        \label{eq:peps_boundary_coupling_inverse}
    \end{align}
\end{theorem}
\begin{proof}\leanok
    Exchange the finite sum over $\nu$ with the product over boundary edges:
    \begin{align}
        \sum_{\nu}K_R^X(\mu,\nu)\,(K_R^X)^{-1}(\nu,\rho)
        &=
        \prod_{f\in\partial R}
        \sum_a G_f^X(\mu_f,a)\,(G_f^X)^{-1}(a,\rho_f).
        \notag
    \end{align}
    Since
    \begin{align}
        \sum_a G_f^X(\mu_f,a)\,(G_f^X)^{-1}(a,\rho_f)
        &=
        (G_f^X(G_f^X)^{-1})_{\mu_f,\rho_f}
        =
        \delta_{\mu_f,\rho_f},
        \notag
    \end{align}
    each factor reduces to the Kronecker symbol at $(\mu_f,\rho_f)$. The
    product of these one-edge Kronecker symbols is $\delta_{\mu,\rho}$.
\end{proof}

\begin{theorem}[Gauge preservation of blocked-region injectivity]
    \label{thm:peps_regionBlockedTensorInjective_applyGauge}
    \lean{TNLean.PEPS.regionBlockedTensorInjective_applyGauge}
    \leanok
    \uses{def:applyGauge, thm:peps_regionBlockedWeight_applyGauge,
        thm:peps_sum_regionBoundaryCoupling_mul_inv}
    Let $\widetilde B=\operatorname{Gauge}_X(B)$. If the blocked tensor family
    $\{T_{B,R}^{\nu}\}_{\nu}$ is linearly independent, then the blocked tensor
    family $\{T_{\widetilde B,R}^{\mu}\}_{\mu}$ is linearly independent.
    Equivalently,
    $\inj(T_{B,R})\Longrightarrow\inj(T_{\widetilde B,R})$.
\end{theorem}
\begin{proof}\leanok
    By (\ref{eq:peps_boundary_coupling}),
    $T_{\widetilde B,R}^{\mu}
    =\sum_{\nu}K_R^X(\mu,\nu)\,T_{B,R}^{\nu}$.
    Suppose $\sum_\mu c_\mu T_{\widetilde B,R}^{\mu}=0$. Substituting
    (\ref{eq:peps_boundary_coupling}) and using the linear independence of
    $\{T_{B,R}^{\nu}\}_{\nu}$ gives
    $\sum_{\mu}c_\mu K_R^X(\mu,\nu)=0$ for every $\nu$.
    Contracting this equality with the right inverse of $K_R^X$ yields, for
    every $\rho$,
    \begin{align}
        c_\rho
        &=
        \sum_{\nu}
        \left(\sum_{\mu}c_\mu K_R^X(\mu,\nu)\right)
        (K_R^X)^{-1}(\nu,\rho)
        =
        0.
        \notag
    \end{align}
    Thus all coefficients $c_\rho$ vanish, proving the linear independence of
    $\{T_{\widetilde B,R}^{\mu}\}_{\mu}$.
\end{proof}

\begin{theorem}[Kernel descent for blocked regions]%
    \label{thm:peps_regionBlockedTensorInjective_of_isVertexInjective}
    \lean{TNLean.PEPS.regionBlockedTensorInjective_of_isVertexInjective}
    \leanok
    \uses{def:IsVertexInjective, def:peps_regionInjectivityDataOf}
    Let $A$ be vertex-injective, and assume that every virtual bond dimension is
    positive. Then every finite region $R$ has an injective blocked tensor:
    $\inj(R)$.
% Paper-gap note: docs/paper-gaps/peps_injective_ft_section3_route.tex.
\end{theorem}
\begin{proof}\leanok
    Let $(c_\rho)_\rho$ be a finitely supported coefficient family on the
    boundary configurations of $R$, and assume that, for every physical
    configuration $\tau$,
    $\sum_{\rho} c_\rho\,T_{A,R}^{\rho}(\tau)=0$.
    The kernel descent removes the vertices of $R$ one at a time and preserves
    the vanishing contraction. After all vertices have been removed, the
    terminal contraction gives $c_\rho=0$ for every boundary configuration
    $\rho$. Hence the family $\{T_{A,R}^{\rho}\}_\rho$ is linearly independent,
    which is $\inj(R)$.
\end{proof}

\begin{theorem}[Union closure for a vertex-injective tensor]%
    \label{thm:peps_regionInjectivityUnionClosure_of_isVertexInjective}
    \lean{TNLean.PEPS.regionInjectivityUnionClosure_of_isVertexInjective}
    \leanok
    \uses{def:IsVertexInjective, def:peps_regionInjectivityDataOf,
        thm:peps_regionBlockedTensorInjective_of_isVertexInjective}
    Let $A$ be vertex-injective with every bond dimension positive. Then the
    concrete region-injectivity predicate of $A$ satisfies union closure: for
    any two finite regions $R$ and $S$,
    $\inj(R)\wedge\inj(S)\Longrightarrow\inj(R\cup S)$.
% Paper-gap note: docs/paper-gaps/peps_injective_ft_section3_route.tex.
\end{theorem}
\begin{proof}\leanok
    The kernel descent theorem gives $\inj(Q)$ for every finite region $Q$.
    Taking $Q=R\cup S$ gives $\inj(R\cup S)$, so the union-closure implication
    holds independently of the injectivity of $R$ and $S$.
\end{proof}

\begin{theorem}[Union closure from the overlapping union lemma]%
    \label{thm:peps_regionInjectivityUnionClosure_of_overlap}
    \lean{TNLean.PEPS.regionInjectivityUnionClosure_of_overlap}
    \leanok
    \uses{def:peps_regionInjectivityDataOf,
        lem:peps_injectiveRegion_union_pos}
    Let $A$ be a PEPS tensor with every bond dimension positive. Then the
    concrete region-injectivity predicate of $A$ satisfies union closure: for
    any two finite regions $R$ and $S$,
    $\inj(R)\wedge\inj(S)\Longrightarrow\inj(R\cup S)$.
    Thus injectivity of the two regions implies injectivity of their union;
    no single-vertex injectivity hypothesis is needed.
    % Source-faithful scope note: this is the union-of-injective-regions
    % lemma, rather than the stronger vertex-injective specialization.
\end{theorem}
\begin{proof}\leanok
    The union-closure implication is the overlapping union lemma
    (Lemma~\ref{lem:peps_injectiveRegion_union_pos}): injectivity of $R$ and of
    $S$, together with the positive bond dimensions, gives injectivity of
    $R\cup S$.
\end{proof}

\begin{lemma}[Finite unions of injective regions]%
    \label{lem:peps_regionInjectivityUnionClosure_finite}
    \lean{TNLean.PEPS.RegionInjectivityUnionClosure.biUnion_injective}
    \leanok
    \uses{def:peps_regionInjectivityData}
    Under the union-closure hypothesis for region injectivity, any nonempty
    finite union of injective finite regions is injective.
\end{lemma}
\begin{proof}\leanok
    Induct on the nonempty finite indexing set. The singleton case is the
    given injectivity hypothesis, and the induction step is one application of
    binary union closure.
\end{proof}

\begin{lemma}[Indexed inside union under closure]%
    \label{lem:peps_regionInjectivityUnionClosure_inside}
    \lean{TNLean.PEPS.RegionInjectivityUnionClosure.regionUnionPart_inside_injective}
    \lean{TNLean.PEPS.RegionInjectivityUnionClosure.regionUnionPart_inside_binary_injective}
    \leanok
    \uses{def:peps_regionInjectivityData,
        lem:peps_regionUnionDecomposition_identities}
    Under the union-closure hypothesis for region injectivity, if the indexed
    subfamilies $\{0,1\}$ and $\{1,2\}$ in the four-region decomposition
    are injective, then the inside subfamily $\{0,1,2\}$ is injective.
    Equivalently, injectivity of the two source-paper regions
    $(A\setminus B)\cup(A\cap B)$ and
    $(A\cap B)\cup(B\setminus A)$ gives injectivity of
    $(A\setminus B)\cup(A\cap B)\cup(B\setminus A)=A\cup B$.
\end{lemma}
\begin{proof}\leanok
    Rewrite the three selected indexed unions as $A$, $B$, and $A\cup B$,
    then apply binary union closure.
\end{proof}
